\section{Methodology}
\label{sec:method}
% EDIT (2026-10-29): new section built from the old Sections 3 (pipeline, profiling data, prompt
% structure), 5.1 (prompt variants), 6.2 (constraint text), 4.3--4.5 (correctness, timing, reporting
% conventions), 7.2--7.3 (gates, application timing) and the "Arms"/"Setup" paragraphs of 8.1 and 8.2,
% following the reviewer's requested structure. The variant descriptions were re-derived from the
% saved prompts (diffs of src/results/atan2-codex-gpt5-<variant>/*prompt_1.txt against the cot prompt)
% and correct two statements of the earlier text: no_constraints removes only the six rules written
% as "Do NOT", and minimal keeps the hardware guidance and the entry-point guardrails.
This section describes the system under study---the optimization loop, the profiling data it can
place in the prompt, and the prompt itself---and then the prompt variants, the correctness checks and
the measurement protocol on which every result depends.

\subsection{The Optimization Loop}
\label{sec:subject}
The framework takes a source file, a build-and-run script and the profile of the program, and closes
the loop shown in Figure~\ref{fig:pipeline}: \emph{profile $\rightarrow$ summarize $\rightarrow$ prompt
$\rightarrow$ generate $\rightarrow$ build $\rightarrow$ verify $\rightarrow$ measure}, with bounded
repair. Profiling happens once, before the loop. The prompt carries the source file, the profile
summaries and the instructions described below; the model answers with a complete replacement for the
file. That file replaces the original in a private copy of the program tree, so the build, the callers
and the input data are those of the real program. The build-and-run script then compiles the program,
runs it, and applies its correctness check. If compilation fails, the program crashes, or the check
prints a failure, the loop summarizes the build log (at most $80$ lines), assembles a repair prompt
from the original source, the rejected candidate and that summary, and asks the model again; a repair
that returns the candidate unchanged is rejected and requested once more. We allow three iterations
before recording the run as a failure. When a candidate builds and passes, the loop records its
in-loop execution time, the average of ten runs of the binary. That number serves the loop only; the
timings reported in this paper come from a separate interleaved measurement pass
(Section~\ref{sec:timing}).

\begin{figure}[t]
\centering
\resizebox{\textwidth}{!}{%% Workflow of the optimization loop (Figure 1). NEW (2026-10-29), drawn from src/opal_cli_codex.py:
%% prompt -> LLM -> candidate -> build+run script (build, run, built-in check or gate) -> on error,
%% summarized log -> repair prompt -> LLM, at most three iterations -> in-loop 10-run average.
%% Profiling (two Nsight Compute passes) runs before the loop; the reported timings come from a
%% separate interleaved pass. Needs \usetikzlibrary{arrows.meta,positioning,calc} (main.tex).
%% Layout: row 1 profile..candidate; row 2 (right to left) build, verify, measure, timing pass;
%% row 3 the repair step under build. No arrow crosses a node; the loop back to the model runs
%% above row 1. NOT compiled yet (no LaTeX here): if a build fails, check this file first.
\begin{tikzpicture}[
  font=\footnotesize,
  box/.style={draw, rounded corners=2pt, align=center, minimum height=2.2em, inner sep=4pt,
              text width=5.6em},
  wide/.style={box, text width=7.6em},
  dec/.style={draw, align=center, minimum height=2.2em, inner sep=3pt, text width=5.2em},
  arr/.style={-{Latex[length=2mm]}, thick},
  lab/.style={font=\scriptsize, fill=white, inner sep=1pt},
  node distance=5mm and 7mm
]
  % row 1: before the loop, and the generation step
  \node[wide] (prof)   {Profile\\(Nsight Compute,\\two passes)};
  \node[wide, right=of prof] (summ) {Summarize\\(PC sampling, roofline,\\kernel importance)};
  \node[wide, right=of summ] (prompt) {Assemble prompt\\(source + summaries\\+ guardrails + format)};
  \node[box,  right=of prompt] (llm) {LLM};
  \node[box,  right=of llm] (cand) {Candidate\\source file};
  % row 2: build, verify, measure, and the separate timing pass
  \node[box,  below=of cand] (build) {Build};
  \node[dec,  left=of build] (verify) {Verify\\(built-in check\\or gate)};
  \node[wide, left=of verify] (measure) {Measure\\(in-loop average\\of 10 runs)};
  \node[wide, left=of measure] (report) {Interleaved timing\\pass (all arms,\\one session)};
  % row 3: repair
  \node[wide, below=of build] (repair) {Summarize log,\\build repair prompt};

  \draw[arr] (prof) -- (summ);
  \draw[arr] (summ) -- (prompt);
  \draw[arr] (prompt) -- (llm);
  \draw[arr] (llm) -- (cand);
  \draw[arr] (cand) -- (build);
  \draw[arr] (build) -- node[lab, above] {ok} (verify);
  \draw[arr] (verify) -- node[lab, above] {pass} (measure);
  \draw[arr] (measure) -- node[lab, above] {reported} (report);
  \draw[arr] (build) -- node[lab, right] {error} (repair);
  \draw[arr] (verify.south) |- node[lab, pos=0.25, left] {fail} (repair.west);
  \draw[arr] (repair.east) -- ++(5mm,0) |- ($(cand.north)+(0,4mm)$)
             -| node[lab, pos=0.25, above] {repair prompt ($\le 3$ iterations)} (llm.north);
\end{tikzpicture}
}
\caption{The optimization loop. Profiling and summarization run once; the prompt, model, build and
  verification steps form the repair loop, bounded at three iterations. The in-loop time is used only
  to record the run; every timing reported in this paper comes from the separate interleaved pass of
  Section~\ref{sec:timing}.}
\label{fig:pipeline}
\end{figure}

\subsection{Profiling Data}
\label{sec:profdata}
% EDIT (2026-10-29): more context on what each kind of data is, asked for in review.
The profile comes from Nsight Compute, NVIDIA's kernel profiler, in two passes over the program: an
uncapped pass that records every kernel launch and gives true launch counts and per-launch durations,
and a capped \texttt{-{}-set full} pass over a bounded number of launches that collects the detailed
metric pages. From the two reports we derive three kinds of prompt context, each rendered as a short
text or JSON summary:
\begin{itemize}\itemsep2pt
  \item \textbf{PC sampling} (\texttt{-{}-page source}). The profiler periodically samples the program
    counter of running warps and records, for each sample, whether the warp was stalled and why (waiting
    on memory, on a math unit, on a barrier, and so on). Aggregated per source line, this says which
    lines the kernel spends its time on and what those lines are waiting for. The summary lists the
    lines with the highest stall ratios, each with its source text and its stall reasons, so that the
    model can refer to them by line number.
  \item \textbf{Roofline} (\texttt{-{}-page details}). For each kernel, the achieved DRAM and compute
    throughput relative to the device's peaks, which says whether the kernel is bound by memory
    bandwidth, by arithmetic throughput, or by neither (latency-bound), together with the analysis rules
    that Nsight Compute itself reports for the kernel.
  \item \textbf{Kernel importance}. The aggregate cost of each kernel (launch count $\times$ per-launch
    time), used to rank which kernel is worth editing at all. For the single-kernel benchmarks this is
    trivial; for the applications it is what selects the target file.
\end{itemize}
The arms that remove or extend this evidence are defined with the prompt variants below.

\subsection{Prompt Structure}
\label{sec:promptstructure}
% EDIT (2026-10-29): the regions are now listed in the order they appear in the prompt, with an excerpt
% of each (Table~\ref{tab:prompt}), taken from the saved prompt
% src/results/atan2-codex-gpt5-cot/atan2_pc+rl_gpt5_cot_prompt_1.txt.
Table~\ref{tab:prompt} shows the full prompt region by region, with the text of each region as sent.
A one-line persona is sent as the system message. The user message then carries, in order: the
numbered source file; the profile summaries; a hardware-guidance paragraph naming the Hopper features
to target; five entry-point guardrails that protect the program's \texttt{main} and its linkage; and
the instruction block. The instruction block states the task and requires a correctness argument
(preconditions, postconditions, loop invariants and a proof that the transformation preserves them),
and then gives nineteen numbered rules: six prohibitions written as ``Do NOT'' (no new top-level
symbols, no wrapper helpers, no larger shared-memory footprint, and three entry-point rules), a set of
working rules (use the profile's line references, make at least one substantive change, keep the code
unchanged and say so if none is justified, no line numbers, no renamed functions or changed
signatures), and the output format: the complete source in one fenced code block followed by two
structured lists, the optimizations applied and those considered but not applied.

\begin{table}[t]
\centering
\footnotesize
\caption{The prompt, region by region, with the text as sent (excerpts; $\langle\cdot\rangle$ marks
  generated content). The persona is the system message; everything else is the user message, in this
  order.}
\label{tab:prompt}
\begin{tabular}{@{}p{0.2\textwidth}p{0.76\textwidth}@{}}
\toprule
\textbf{Region} & \textbf{As sent} \\
\midrule
Persona (system message) & \texttt{You are an HPC performance optimization expert.} \\
Source file & \texttt{\# Begin Source Code} $\langle$the file, one numbered line per source line$\rangle$
  \texttt{\# End Source Code} \\
Profile summaries & \texttt{\# Begin Performance Analysis} $\langle$PC-sampling summary: for each hot
  line, its number, its source text and its stall reasons$\rangle$ $\langle$roofline summary$\rangle$
  \texttt{\# End Performance Analysis} \\
Hardware guidance & \texttt{Optimize specifically for NVIDIA H100 (Hopper) by targeting L2/cache reuse,
  HBM3 bandwidth, and Tensor Core throughput. Prioritize tiling and data locality: \ldots} \\
Entry-point guardrails & \texttt{1. The source already contains 1 main() definition(s). Preserve exactly
  that many main() definitions. 2. Do NOT add a new main(), duplicate main(), rename main(), \ldots} \\
Task and proof requirement & \texttt{Optimize the following HPC code for execution time. You must define
  mathematical preconditions and postconditions, provide explicit loop invariants to guarantee semantic
  equivalence, and provide a logical proof showing why the invariants remain unbroken \ldots} \\
Prohibitions & \texttt{9. Do NOT introduce any new helper functions, utility functions, macros,
  templates, classes, structs, global variables, or other new top-level symbols \ldots\
  11. Do NOT increase the total shared-memory footprint of any kernel \ldots} \\
Working rules & \texttt{12. Use the performance analysis reference number(s) to guide your changes.
  \ldots\ 15. When the performance data supports it, apply at least one substantive
  performance-motivated code change \ldots\ 16. If no safe substantive optimization can be justified
  \ldots, keep the code semantically unchanged and say so explicitly \ldots} \\
Output format and explanation & \texttt{17. Return: (a) the full source code including the optimized
  kernels inside {\textasciigrave\textasciigrave\textasciigrave} block (b) optimizations list = [\{'lines': [...], 'reason': '...'\}] (c)
  suggested\_but\_not\_applied list = [\{'lines': [...], 'reason': '...'\}]} \\
Step-by-step instruction (\texttt{cot} only) & \texttt{Think step by step about which bottlenecks to
  address and why before writing any code.} \\
\bottomrule
\end{tabular}
\end{table}

\subsection{Prompt Variants}
\label{sec:ablation}
To find out which parts of the prompt matter, the prompt-component study evaluates six variants, each
of which removes or adds exactly one region while leaving the source file and the profile summaries in
place. Table~\ref{tab:variants} shows which regions each variant carries. We describe them from the
least populated prompt to the most populated one; the effect of each is reported in
Section~\ref{sec:prompts}.

\begin{table}[t]
\centering
\footnotesize
\setlength{\tabcolsep}{4pt}
\caption{Prompt regions carried by each variant of the prompt-component study, from the least to the
  most populated. The persona, the source file, the profile summaries and the single-code-block output
  requirement are present in every variant. \cmark\ = present, \xmark\ = absent.}
\label{tab:variants}
\begin{tabular}{@{}lcccccc@{}}
\toprule
\textbf{Region} & \texttt{minimal} & \texttt{no\_role} & \texttt{no\_constraints} & \texttt{no\_proof}
  & \texttt{no\_explanation} & \texttt{cot} \\
\midrule
Hardware guidance            & \cmark & \xmark & \cmark & \cmark & \cmark & \cmark \\
Entry-point guardrails       & \cmark & \cmark & 2 of 5 & \cmark & \cmark & \cmark \\
Proof requirement            & \xmark & \cmark & \cmark & \xmark & \cmark & \cmark \\
Prohibitions (``Do NOT'')    & \xmark & \cmark & \xmark & \cmark & \cmark & \cmark \\
Working rules                & \xmark & \cmark & \cmark & \cmark & \cmark & \cmark \\
Explanation lists            & \xmark & \cmark & \cmark & \cmark & \xmark & \cmark \\
Step-by-step instruction     & \xmark & \xmark & \xmark & \xmark & \xmark & \cmark \\
\bottomrule
\end{tabular}
\end{table}

\paragraph{Minimal (\texttt{minimal}).}
The whole instruction block---task statement, proof requirement, the nineteen rules and the output
format---is replaced by one sentence: \textit{``Optimize the following CUDA code for execution time.
Return the complete optimized source in a single \texttt{cpp} block.''} The source, the profile
summaries, the hardware guidance and the entry-point guardrails remain, so \texttt{minimal} isolates
the value of the \emph{instructions}; removing the evidence instead is the separate profile-data
ablation (\texttt{code\_only}, below).

\paragraph{No Role (\texttt{no\_role}).}
Removes the hardware-guidance paragraph, which tells the model which Hopper features to target (cache
reuse, asynchronous copies, Tensor Cores, occupancy, warp efficiency). The one-line persona is a system
message and stays in every variant.

\paragraph{No Constraints (\texttt{no\_constraints}).}
Removes the six rules written as prohibitions: no new helper functions, macros, templates or other
top-level symbols; no wrapper load helpers; no larger shared-memory footprint; no second entry point;
no wrapper executables; and no renaming of externally referenced functions. The remaining rules
(including ``do not change any function signatures'' and ``do not rename or duplicate any kernel'')
stay, so the variant frees the model to add code around the existing kernels rather than to change
their interface.

\paragraph{No Proof (\texttt{no\_proof}).}
Removes the sentence that requires preconditions, postconditions, loop invariants and a proof of
semantic equivalence. The model may then apply a transformation without having to argue that it
preserves the computation.

\paragraph{No Explanation (\texttt{no\_explanation}).}
Removes the two structured lists from the required output---the optimizations applied and those
considered but not applied---and asks for the code alone. Having to name and justify each change is a
form of implicit reasoning, and the lists are also the record from which we reconstruct \emph{why} a
run behaved as it did.

\paragraph{Chain of Thought (\texttt{cot}).}
Unlike the others, \texttt{cot} adds a region: the instruction \textit{``Think step by step about which
bottlenecks to address and why before writing any code''} is prepended to the instruction block. It is
the most populated prompt in the study and, since the study has no separate full-prompt arm, the
closest thing to one.

\paragraph{Variants not evaluated.}
Two further variants are implemented but excluded from the measured grid. \texttt{no\_format} strips
the output-format requirement and asks only for optimized code; it perturbs the reliability of
automated code extraction rather than the quality of the optimization, which would confound the metric.
\texttt{no\_hardware} removes the hardware-guidance paragraph and the target-architecture line; in the
prompt format used here it removes the same paragraph as \texttt{no\_role}, so it is subsumed.

\paragraph{Variants used in the later experiments.}
Three prompts defined after the component study are used in Sections~\ref{sec:realapps}
to~\ref{sec:guidvar}:
\begin{itemize}\itemsep2pt
  \item \textbf{Baseline} is the full prompt of Table~\ref{tab:prompt} with PC-sampling and roofline
    summaries (\texttt{pcrl} in the profile-data ablation, where the name distinguishes it from the arm
    without data).
  \item \textbf{Structural} relaxes the prohibitions: it allows new kernels, helpers and templates,
    signature and launch-configuration changes, and shared-memory staging, under an explicit contract
    that the computation is preserved, and it adds two derived summaries, a per-kernel budget (launch
    count, aggregate time, DRAM and SM utilization, bound verdict) and an occupancy summary (achieved
    occupancy, waves per SM, and the resource that limits occupancy).
  \item \textbf{Diagnostics} is the structural prompt plus a bounded checklist of measured conditions
    and the transformation each one suggests. Each item is triggered by the kernel's own numbers
    (occupancy, DRAM and SM throughput, shared-memory use) and names a remedy, for example
    \textit{``Is a value read from and written back to global memory inside the innermost loop (an
    accumulator, counter or running sum)? \ldots\ Accumulate in a local variable and store once after
    the loop.''} It is our instance of profiling data made prescriptive.
    % Quoted from src/DiagnosticsSummarization.py, item (a).
  \item \textbf{Code only} (\texttt{code\_only}) is the baseline prompt without the profile summaries.
    Removing the summaries alone would bias the comparison, because four instruction sentences refer to
    the data---most importantly rule 16, \textit{``if no safe substantive optimization can be
    justified from the provided performance data, keep the code semantically unchanged''}, which would
    invite a no-op when no data is present. Those four sentences are reworded to refer to the model's
    own analysis of the source; the two prompts are otherwise byte-identical (verified: the diff is the
    removed block plus those four sentences).
\end{itemize}

\subsection{The Semantic-Preservation Constraint}
\label{sec:constraint}
% EDIT (2026-10-29): constraint text moved here from the old Section 6.2; the dates of the arms that
% carry it are from the saved prompts on the cluster (bincount-gpt5 2026-09-09: absent;
% bincount-gpt5-struct 2026-09-09, autodock_gpu-gpt5 2026-09-17, *-prof-* 2026-09-22: present).
The gaming described in the introduction led us to add a \emph{semantic preservation} block to the
prompt. It states that the model may change \emph{how} a computation executes (data placement, memory
layout and access pattern, thread/block mapping, vectorization, unrolling, loop fusion, tiling and
reordering, hoisting loop-invariant work) but may not change \emph{what} is computed or how much work
is done. In particular it may not reduce the number of algorithm-level operations by exploiting
algebraic identities, change the iteration count of a timing, repeat or convergence loop, precompute at
compile time a value the program computes at run time, skip work on the grounds that the program's own
verification would still pass, or alter printed output. A closing sentence names the case directly: if
a benchmark exists to measure a specific hardware operation, that operation must still be performed the
same number of times on the same volume of data. The rule is deliberately scoped to \emph{benchmark
scaffolding}; eliminating genuine redundancy inside an application remains allowed, since that is the
goal for real codes (Section~\ref{sec:gaming-disc}). The block was added after the component study.
The six variants of that study and the baseline arms generated before the change do not carry it;
the structural, diagnostics, \texttt{pcrl} and \texttt{code\_only} prompts and the application study
all do.

\subsection{Correctness Validation}
\label{sec:correctness}
% EDIT (2026-10-29): merged from the old 4.3 and 7.2.
A candidate counts only if it is correct, and correctness is checked in whichever way the program
supports. The HeCBench benchmarks ship a built-in check that prints PASS or FAIL after comparing the
GPU result with a host reference, or print numerical results that the candidate must reproduce; the
build-and-run script applies that check in every iteration of the loop, so a failing candidate is sent
back for repair rather than timed. Benchmarks whose check cannot pass for the unmodified program are
excluded (Section~\ref{sec:guidvar}).

Applications do not ship a check, and their outputs are not always reproducible, so we built one gate
per application and validated each gate negatively as well as positively:
\begin{itemize}\itemsep2pt
  \item \textbf{AutoDock-GPU} is stochastic on the GPU even with a fixed seed (best energies
    $-13.89$, $-14.38$, $-14.01$\,kcal/mol for one binary), so the gate is a band: 50 poses and best
    energy $\le -12.5$\,kcal/mol. Negative test: flipping the gradient sign yields $-0.53$ (FAIL).
  \item \textbf{Cholla} is bit-deterministic; the gate compares the full 3D snapshot field by field
    (relative $L_1 \le 10^{-9}$, relative max $\le 10^{-6}$) after 137 steps. Negative tests: a flipped
    energy-flux sign diverges (FAIL) and a $10^{-6}$ relative perturbation of one flux is caught
    ($L_1 \approx 10^{-6}$, three orders of magnitude over tolerance).
\end{itemize}
Every application candidate reported in this paper passed its gate. Since a gate checks the outputs the
program prints and not all internal state, ``gated correct'' means ``passed these checks'', not
``verified''; Section~\ref{sec:limits} discusses what that leaves open.

\subsection{Timing and Measurement Protocol}
\label{sec:timing}
% EDIT (2026-10-29): merged from the old 4.4, 4.5 and 7.3, stated once; "drift" defined before it is
% used and the 49% figure clarified (the same unmodified binary timed in two sessions, HeCBench study,
% scripts/optimas_study/README.md item 5). "Pedantic" removed.
\paragraph{What is timed.}
Time comes from the program's own reported kernel or phase time wherever it provides one, and from wall
clock otherwise. For the applications this matters: AutoDock-GPU's CUDA setup alone takes $68$\,s on a
cold driver and would swamp a few percent of solver time, so we use each application's own phase timer.

\paragraph{Drift, and the protocol that controls it.}
By \emph{drift} we mean a change in the measured time of the same binary between two sessions, with no
change to the code. Our GPU nodes are shared with other users, and drift there is large: the same
unmodified HeCBench binary timed in two sessions differed by up to $49\%$, an order of magnitude more
than most of the effects we are trying to measure. Two things follow. First, the pipeline's in-loop
time is never reported; it was off by $15\%$ for one application arm. Second, every reported number,
except the prompt-component grid of Section~\ref{sec:prompts} (whose runtimes are the in-loop
ten-run averages, measured during that study), comes from a separate measurement pass with one design:
build every arm of a program once; discard a warm-up run; wait for an idle GPU before starting; time
all arms round-robin within a single session, so that whatever the node is doing affects every arm
alike; log other users' GPU processes at the start \emph{and} the end of every repetition, and exclude
the repetitions that were not clean; and take the median over the clean repetitions. Repetitions that
disagree with their neighbours are treated as suspect even when no foreign process was seen. The number
of clean repetitions behind each median is stated with each table: five per arm for the benchmarks, and
for the applications between two and seven, depending on how often other users' jobs intervened. Two
complete timing passes were discarded because another user's job occupied the GPU throughout.

\paragraph{Reporting conventions.}
Speedups are reported against the unmodified program unless a table states otherwise. Except where
stated, each (program, model, prompt) cell is a single language-model sample, and an effect below
roughly $5\%$ in a single cell is treated as unresolved rather than small. Where two arms are compared
on a noisy benchmark, the pair is scored a \emph{tie} when the difference between the arms,
$|\mathrm{a}/\mathrm{b}-1|$, is smaller than the largest repetition spread measured for that pair,
that is, when the difference does not exceed its own measurement noise. Finally, any speedup above
$5\times$ is treated as a correctness alarm rather than a result: its diff is read before it is
reported, and Section~\ref{sec:gaming-disc} explains why.
